\section{Competition between pools with dividend strategy}

\subsection{Wealth processes with diversification}
\subsubsection{miners' wealth process}

Consider $m \in \mathbb{N}$ mining pools, each offering a PPS contract characterized by a difficulty reduction $\delta_k$ and a fee $f_k$, for $k = 1, \dots, m$. Consider also a population of $n$ miners, indexed by $i = 1, \dots, n$. The wealth process \eqref{eq:surplus_solo_in_pool_mining} can be viewed as particular cases of a more general framework in which the miner allocates his computational resources $\lambda_i$ across multiple pools ($k = 1,\cdots,m$) and solo ($k=0$) simultaneously, according to a weight vector
\[
w = (w_{i,0}, w_{i,1}, \dots, w_{i,m}) \in \Delta_{i,m},
\]
where $\Delta_{i,m} = \left\{ w \in \mathbb{R_+}^{m+1} : \sum_{k=0}^m w_{i,k} = 1 \right\}$ denotes the unit simplex. The reward arrival rate of miner $i$ associated with pool $k$ is then given by
\[
\lambda_{i,k} = \frac{\lambda_i w_{i,k}}{\delta_k}, \quad k = 0, \dots, m.
\]

The $\lambda_{i,k}$ define the intensities of conditionally independent Poisson processes given the allocation $w$. Aggregating these contributions, the total reward arrival process of the miner $(\hat{N}^{(i)}_t)_{t \geq 0}$ is a Poisson process with intensity
\[
\hat{\lambda}_i = \sum_{k=0}^m \lambda_{i,k}.
\]
Each reward originates from pool $k$ with probability $p_{i,k} = \lambda_{i,k} / \hat{\lambda} _i$ and has magnitude $b_k = b(1-f_k)\delta_k$. Therefore, the miner’s wealth process can be written as
\begin{equation}\label{welath_multi_miner}
X_t = x - c \cdot t + \sum_{j=1}^{\hat{N}^{(i)}_t} B^{(i)}_j,
\end{equation}
where $(B_j)_{j \geq 1}$ is a sequence of independent and identically distributed random variables satisfying
\[
\mathbb{P}(B^{(i)} = b_k) = \frac{\lambda_{i,k}} {\hat{\lambda}_i} = p_{i,k}, \quad k = 0, \dots, m.
\]

Figure \ref{fig:multipool_structure} illustrates how the wealth process evolves in this architecture.

\begin{figure}[htbp]
    \centering
    \includegraphics[width=0.7\textwidth]{multipool_diagram.png}
    \caption{Multi-pool mining architecture.}
    \label{fig:multipool_structure}
\end{figure}

\subsubsection{Pool managers' wealth process}

Under this allocation, the total hashpower received by pool $k$ is given by
\[
H_k(w)=\sum_{i=1}^n \lambda_i w_{i,k}.
\]
As in the single-pool participation setting, the arrival rate of shares in pool $k$ is determined by the effective difficulty parameter $\delta_k$, yielding a Poisson intensity
\[
\mu_k(w)=\frac{H_k(w)}{\delta_k}
=\frac{1}{\delta_k}\sum_{i=1}^n \lambda_i w_{i,k}.
\]

Each share submitted to pool $k$ is either a valid block with probability $\delta_k/(1+\delta_k)$, or a non-block share with probability $1/(1+\delta_k)$. Regardless of its nature, the pool manager pays a fixed PPS reward $b_k=b\,\delta_k(1-f_k)$ to the contributing miner. In the event of a valid block, the manager additionally receives the block reward $b$. The wealth process of the manager of pool $k$ is therefore given by the compound Poisson process
\begin{equation}\label{eq:surplus_pool_manager_split}
Y_t^{(k)} = y_k + \sum_{j=1}^{M_t^{(k)}} V_j^{(k)},
\end{equation}
where, conditionally on the allocation profile $w$,
\begin{itemize}
    \item $(M_t^{(k)})_{t\geq 0}$ is a Poisson process with intensity $\mu_k(w)$,
    \item the sequence $(V_j^{(k)})_{j\geq 1}$ is i.i.d. with distribution
    \[
    \mathbb{P}\!\left(V^{(k)}=-b_k\right)=\frac{1}{1+\delta_k},
    \qquad
    \mathbb{P}\!\left(V^{(k)}=b-b_k\right)=\frac{\delta_k}{1+\delta_k}.
    \]
\end{itemize}

We emphasize that, in contrast to the single-pool participation case, the total hashpower $H_k(w)$ is no longer determined by a partition of miners, but rather by the aggregation of fractional contributions. This induces a coupling across pools through the allocation profile $w$, while preserving conditional independence of the underlying Poisson processes.



\subsection{Expected discounted dividend maximization}
\subsubsection{miners' dividend strategy}

A natural performance criterion for the miner is given by the classical dividend optimization problem from risk theory. The wealth process \eqref{welath_multi_miner} can be interpreted as the surplus of the miner, from which cash flows are extracted over time in the form of dividends. Formally, a dividend strategy is described by a non-decreasing, adapted process $(L_t)_{t \geq 0}$, representing the cumulative payments made up to time $t$. The controlled surplus process then writes
\[
    U_t = X_t - L_t,
\]
and ruin occurs at the stopping time
\[
    \tau = \inf\{t \geq 0 : U_t < 0\}.
\]
The objective of the miner is to maximize the expected discounted dividends
until ruin, namely
\begin{equation}\label{eq:dividend_value}
    V(x) = \sup_{L} \mathbb{E}_x\left[\int_0^{\tau} e^{-rt} \, dL_t \right],
\end{equation}
where $r \geq 0$ is a discount rate. For any fixed allocation profile $w \in \Delta_{i,m}$, the process $(X_t)_{t \geq 0}$ is a spectrally positive L\'{e}vy process. In this setting, it is well established that the optimal dividend strategy is of barrier type (see \citet{avanzi2007optimal} and \citet{bayraktar2013optimal}): there exists a level $a_w^* > 0$ such that it is optimal to pay out immediately any surplus above $a_w^*$, while no dividends are paid when the surplus lies below this level. The controlled process is thus reflected at the barrier $a_w^*$.

Denoting by $V_w(x)$ the value function associated with \eqref{eq:dividend_value} under the optimal barrier strategy, the miner's problem reduces to a static optimization over the allocation set, since each $w \in \Delta_{i,m}$ uniquely determines an optimal barrier $a_w^*$ and a corresponding value $V_w(x)$. The optimal allocation is therefore
\[
    w^* = \arg\max_{w \in \Delta_{i,m}} V_w(x).
\]
The contribution of \citet{goffard2025hashpower} is to show that this problem can be tractably analyzed in the context of PPS mining pools. By exploiting the L\'{e}vy structure of the wealth process, the value function $V_w(x)$ can be characterized in terms of scale functions. This representation allows for the computation of both the optimal barrier level $a_w^*$ and the associated value $V_w(x)$, thereby enabling a systematic comparison of PPS contracts and a characterization of the optimal allocation $w^*$.

\subsubsection{pool managers' dividend strategy}

Having characterized the wealth process \eqref{eq:surplus_pool_manager_split} of the pool
manager, we now turn to the associated dividend optimization problem. The surplus process of
pool $k$ is a pure jump process with two-sided increments: upward jumps of size $b - b_k$ occur
upon block discovery, while downward jumps of size $b_k$ occur upon submission of a non-block
share. Since the process is constant between consecutive jump times, no strategic decision needs
to be made in continuous time. It is therefore sufficient to inspect the process at jump times only. Let $(T_i^{(k)})_{i \geq 1}$ denote the jump times of $(M_t^{(k)})_{t \geq 0}$, which form a
sequence of i.i.d.\ exponential random variables with parameter $\mu_k(w)$. The embedded
discrete-time process
\[
    S_n^{(k)} := Y_{T_n^{(k)}}^{(k)} = y_k + \sum_{i=1}^n V_i^{(k)}, \quad n \geq 0,
\]
is a standard random walk. A dividend strategy is then a sequence of decisions
$\mathcal{S} = \{f_1, f_2, \ldots\}$, where each $f_n : \mathbb{R} \to \mathbb{R}_+$ specifies
the dividend paid at step $n$, leading to the controlled process
\[
    R_n^{(k)} = S_{n-1}^{(k)} - f_{n-1}\!\left(S_{n-1}^{(k)}\right) + V_n^{(k)},
    \quad n \geq 1.
\]
Ruin can only occur at jump times, and is defined as
\[
    \tau^{(k)} = \inf\!\left\{n \geq 1 : R_n^{(k)} < 0\right\}.
\]
The objective of the pool manager is to maximize the expected discounted dividends
\[
    V^{(k)}(y_k) = \mathbb{E}_{y_k}\!\left[
        \sum_{n=1}^{\tau^{(k)}} e^{-r T_n^{(k)}} f_n\!\left(S_{n-1}^{(k)}\right)
    \right],
\]
where $r \geq 0$ is the discount rate. This framework coincides with that of
\citet{albrecher2011randomized}, in which the optimal strategy is shown to be a
\textit{band strategy}: dividends are paid according to a state-dependent rule determined by a
finite collection of intervals. Band strategies are, however, difficult to handle analytically. To obtain a more tractable structure, we impose the following simplifying assumptions, which are natural in the context of
a homogeneous PPS pool.

\begin{hp}\label{ass:integer_state}
The state space of $(S_n^{(k)})_{n \geq 0}$ and $(R_n^{(k)})_{n \geq 0}$ is $\mathbb{Z}$.
\end{hp}

Assumption~\ref{ass:integer_state} requires that $b_k \in \mathbb{N}$, that the initial surplus
$y_k$ is proportional to $b_k$, and that $b$ is proportional to $b_k$. Under this assumption, $(S_n^{(k)})_{n \geq 0}$ is an integer-valued random walk with increments taking
values in $\{-1,\, b-b_k\}$ using the taxonomy of \citet{gerber1972games} for which a barrier strategy is optimal as shown in \citet{miyasawa1961economic}: there
exists a threshold $a_k^* \in \mathbb{N}$ such that the pool manager pays out all surplus above
$a_k^*$ as dividends, while retaining the full surplus when it lies below this level. We now characterize explicitly the value function and the optimal barrier level.

Define the vector of monomials
\begin{equation}\label{eq:rho_vector}
    \bm{\rho}^{[y]} := \left(\rho_1^y,\, \rho_2^y,\, \dots,\, \rho_k^y\right)^\top \in \mathbb{C}^k,
\end{equation}
where $\rho_1, \dots, \rho_k \in \mathbb{C}$ denote the $k$ roots of the characteristic equation
\begin{equation}\label{eq:characteristic}
    (r + \mu)\,\rho \;-\; \mu p\,\rho^{k} \;-\; \mu q \;=\; 0,
\end{equation}
with $p = \delta/(1+\delta)$, $q = 1/(1+\delta)$, and $\mu = \mu_k(w)$ the Poisson intensity
defined in \eqref{eq:surplus_pool_manager_split}.

\begin{prop}\label{prop:optimal_barrier}
Let $(\widetilde Y_t^{(k)})_{t \geq 0}$ denote the normalized surplus process of pool $k$ defined by
\[
\widetilde Y_t^{(k)} := \frac{Y_t^{(k)}}{b_k}.
\]
Under Assumption~\ref{ass:integer_state}, it is an integer-valued compound Poisson process with increments in $\{-1,\, k-1\}$, where $k = b/b_k \in \mathbb{N}_{\geq 2}$. Assume
moreover that the drift condition
\begin{equation}\label{eq:drift_condition}
    p(k-1) - q > 0
\end{equation}
holds, ensuring that the process is not absorbed at ruin almost surely. For a fixed
candidate barrier $a \in \mathbb{N}$ with $a \geq k - 1$, the value function under
the barrier strategy at level $a$ takes the form
\begin{equation}\label{eq:value_function_form}
    V^{(k)}(y,\, a) \;=\; \bm{\rho}^{[y]} \cdot \mathbf{c}(a), \quad y = 0, 1, \dots, a,
\end{equation}
where the coefficient vector $\mathbf{c}(a) \in \mathbb{C}^k$ is the unique solution of the
linear system
\begin{equation}\label{eq:linear_system}
    \mathbf{A}(a)\,\mathbf{c}(a) \;=\; \mathbf{d}(a).
\end{equation}
The matrix $\mathbf{A}(a) \in \mathbb{C}^{k \times k}$ and the vector
$\mathbf{d}(a) \in \mathbb{C}^{k}$ are defined row-wise by
\begin{equation}\label{eq:A_matrix}
    \mathbf{A}(a)_j \;=\;
    \begin{cases}
        (r+\mu)\,\bm{\rho}^{[y_j]} - \mu p\,\bm{\rho}^{[y_j + k - 1]},
        & j = 1, \\[0.6em]
        (r+\mu)\,\bm{\rho}^{[y_j]} - \mu p\,\bm{\rho}^{[a]}
        - \mu q\,\bm{\rho}^{[y_j - 1]},
        & j = 2, \dots, k,
    \end{cases}
\end{equation}
\begin{equation}\label{eq:d_vector}
    \mathbf{d}(a)_j \;=\;
    \begin{cases}
        0,
        & j = 1, \\[0.6em]
        \mu p\,(y_j + k - 1 - a),
        & j = 2, \dots, k,
    \end{cases}
\end{equation}
where $y_1 = 0$ and $y_j = a - (k-1) + (j-1)$ for $j = 2, \dots, k$.
The optimal barrier is then given by
\begin{equation}\label{eq:optimal_barrier}
    a_k^* \;=\; \operatorname*{arg\,max}_{a \,\in\, \mathbb{N},\; a \,\geq\, k-1}
    V^{(k)}(y,\, a).
\end{equation}
\end{prop}

\begin{proof}
We derive the value function by analyzing the dynamic programming equation satisfied
by $V^{(k)}(\cdot, a)$ over an infinitesimal time interval $dt$. Since the surplus
process is constant between jump times, the only events to consider are the arrival
or non-arrival of a jump during $[0, dt)$.

\medskip
\noindent\textbf{Step 1: Dynamic programming equation.}
With probability $1 - \mu\,dt + o(dt)$, no jump occurs and the value function is
discounted at rate $r$. With probability $\mu p\,dt + o(dt)$, an upward jump of
size $k-1$ occurs, and with probability $\mu q\,dt + o(dt)$, a downward jump of
size $1$ occurs. Letting $dt \to 0$ and rearranging, one obtains the
difference equation
\begin{equation}\label{eq:diff_eq}
    (r + \mu)\,V^{(k)}(y,a)
    \;=\;
    \mu p\,V^{(k)}(y + k - 1,\, a)
    \;+\;
    \mu q\,V^{(k)}(y - 1,\, a),
    \quad y = 1, \dots, a-(k-1),
\end{equation}
which holds in the interior of the state space, away from the boundaries.
This is a linear homogeneous difference equation of order $k$, whose general
solution is of the form
\begin{equation}\label{eq:general_solution}
    V^{(k)}(y, a) \;=\; \sum_{i=1}^{k} c_i(a)\,\rho_i^{\,y}
    \;=\; \bm{\rho}^{[y]} \cdot \mathbf{c}(a),
\end{equation}
where $\rho_1, \dots, \rho_k$ are the roots of the characteristic
equation~\eqref{eq:characteristic} and $\mathbf{c}(a) = (c_1(a), \dots, c_k(a))^\top$
are coefficients to be determined by boundary conditions (see e.g. \citet{jerri2013linear}).

\medskip
\noindent\textbf{Step 2: Boundary conditions at $y = 0$.}
For $y = 0$, a downward jump leads to immediate ruin, so that
$V^{(k)}(-1, a) = 0$. The difference equation~\eqref{eq:diff_eq} then
simplifies to
\begin{equation}\label{eq:bc_ruin}
    (r + \mu)\,V^{(k)}(0,\, a) \;=\; \mu p\,V^{(k)}(k-1,\, a).
\end{equation}
Substituting the general solution~\eqref{eq:general_solution} into~\eqref{eq:bc_ruin}
yields the first row of the linear system, namely
\begin{equation*}
    \left((r+\mu)\,\bm{\rho}^{[0]} - \mu p\,\bm{\rho}^{[k-1]}\right)
    \cdot \mathbf{c}(a) \;=\; 0,
\end{equation*}
which corresponds to row $j = 1$ of~\eqref{eq:A_matrix}--\eqref{eq:d_vector}
with $y_1 = 0$.

\medskip
\noindent\textbf{Step 3: Boundary conditions near the barrier.}
For $y \in \{a-(k-1)+1, \dots, a\}$, an upward jump of size $k-1$ brings the
surplus above the barrier. Under the barrier strategy, the excess is immediately
paid as a lump-sum dividend and the surplus is reset to $a$, so that
\begin{equation*}
    V^{(k)}(y + k - 1,\, a)
    \;\longrightarrow\;
    V^{(k)}(a,\, a) + (y + k - 1 - a).
\end{equation*}
Substituting into~\eqref{eq:diff_eq} gives, for each such $y$,
\begin{equation}\label{eq:bc_barrier}
    (r+\mu)\,V^{(k)}(y,\,a)
    \;=\;
    \mu p\!\left(V^{(k)}(a,\,a) + y + k - 1 - a\right)
    + \mu q\,V^{(k)}(y-1,\,a).
\end{equation}
Substituting the general solution~\eqref{eq:general_solution}
into~\eqref{eq:bc_barrier} and rearranging yields, for $j = 2, \dots, k$,
\begin{equation*}
    \left(
        (r+\mu)\,\bm{\rho}^{[y_j]}
        - \mu p\,\bm{\rho}^{[a]}
        - \mu q\,\bm{\rho}^{[y_j-1]}
    \right) \cdot \mathbf{c}(a)
    \;=\;
    \mu p\,(y_j + k - 1 - a),
\end{equation*}
with $y_j = a-(k-1)+(j-1)$, which corresponds to rows $j = 2, \dots, k$
of~\eqref{eq:A_matrix}--\eqref{eq:d_vector}.

\medskip
\noindent\textbf{Step 4: Determination of the optimal barrier.}
Combining Steps 2 and 3 yields the complete $k \times k$ linear
system~\eqref{eq:linear_system}, from which the coefficient vector $\mathbf{c}(a)$
is uniquely determined for each candidate barrier $a \geq k-1$.
The corresponding value function $V^{(k)}(y, a)$ is then evaluated via
\eqref{eq:value_function_form}. The optimal barrier $a_k^*$ is defined
as in~\eqref{eq:optimal_barrier}, and is obtained numerically by solving
the system over a finite grid of candidate values and selecting the maximizer.
\end{proof}

\begin{ex}\label{ex:barrier}
We illustrate numerically the characterization of the optimal dividend barrier
for a PPS pool manager in the lattice setting of
Proposition~\ref{prop:optimal_barrier}.

The parameters are set as
\[
b = 40,
\qquad
f = 0.4,
\qquad
k = 10,
\]
\[
\mu = 6,
\qquad
r = 0.05,
\qquad
x_0 = 20.
\]

The difficulty parameter is chosen so that the surplus process evolves on the
lattice $b_k \mathbb{Z}$, namely
\[
\delta = \frac{1}{k(1-f)},
\]
which implies $b_k = b/k = 4$ and normalized increments in $\{-1,\,k-1\}$.

For each candidate barrier level $a$, the value function $V(x_0,a)$ is computed
analytically by solving the linear system associated with the boundary conditions
of Proposition~\ref{prop:optimal_barrier}. The optimal barrier is obtained by
maximizing $V(x_0,a)$ over a finite grid, yielding
\[
a^* = 108,
\qquad
V(x_0,a^*) \approx 77.3.
\]

To validate this result, we estimate the value function via Monte Carlo simulation of the controlled surplus process under the barrier strategy. The jump times are generated according to a Poisson process of intensity $\mu$, and discounted dividends are accumulated until ruin. Confidence intervals at the $95\%$ level are computed for each barrier. Figure~\ref{fig:barrier_validation} compares the analytical value function with the Monte Carlo estimates.

\begin{figure}[H]
\centering
\includegraphics[width=0.6\textwidth]{div_analytic.png}
\caption{Analytical value function $V(x_0,a)$ and Monte Carlo 95\% confidence intervals.}
\label{fig:barrier_validation}
\end{figure}
\end{ex}


